\scanpage{13}
% archive sheet number: 12
5. \uline{Classe fondamentale locale}.

Reprenons l'homomorphisme de Gysin général (4.3), et \unsure{notans|notons} que par la formule d'adjonction, sa connaissance équivaut à celle d'un homomorphisme
\[
f^{\handwritten{*}}(G)(d-d')\handwritten{[}2(d-d')\handwritten{]}\handwritten{\unsure{\Longrightarrow|\longrightarrow}}\mathrm{R}f^{\handwritten{!}}G,
\tag*{(4.4)}
\]
\margin{\handwritten{\tikz[baseline=(n.base)]{\node[rotate=12] (n) {autre nom?};\draw[rotate=12] (0,0) ellipse (1.15cm and .35cm);}}}
\deleted{\unsure{qui, pour $f$}}\handwritten{[}On l'appelle l'homomorphisme \uline{classe fondamentale locale}(?)\handwritten{]}. Il est fonctoriel en $G$, et, comme nous avons vu au par.4, il suffit de le connaitre pour $G=A_Y$ pour le connaitre pour tout $G$. Dans ce cas, on peut l'interpréter comme un homomorphisme
\[
\handwritten{\phi}_f:\quad A_X\handwritten{\longrightarrow}\mathrm{R}f^{\handwritten{!}}(A_Y)(d'-d)\ 2(d'-d)\handwritten{,}\quad\deleted{\text{\unsure{sur}}\ \unsure{\phi_Y|Y}}
\tag*{(4.5)}
\]
\inserted{\handwritten{i.e. }$\handwritten{\varphi_f\in}\mathrm{H}^{\handwritten{2}(d-d')}(Y,\mathrm{R}f^{\handwritten{!}}(A_Y)(d-d'))$}
et il définit par suite une section, notée de même
\[
\begin{aligned}
&\deleted{\underline{\handwritten{\phi}}_f\ \handwritten{\unsure{\leftarrow|\in}}\ \mathrm{H}^0(Y,\underline{\mathrm{H}}^0(\mathrm{R}f^{\handwritten{!}}(A_Y)(d'-d)\ 2(d'-d))}\quad\\
&\inserted{\handwritten{\varphi_f\in}\mathrm{H}^0(Y,\underline{\mathbb{R}}^{\handwritten{2}(d'-d)}(\deleted{\mathrm{R}}f^{\handwritten{!}}(A_Y))(d'-d))},
\end{aligned}
\tag*{\tikz[baseline=(n.base)]{\node (n) {(4.6)};\handwritten{\draw[gray] (-.65,-.85) .. controls (-.42,-.9) and (-.8,.15) .. (-.6,.3) .. controls (-.25,.65) and (.45,.6) .. (.65,.3) .. controls (.72,.12) and (.55,-.55) .. (.55,-.77) -- (.73,-.81);}}}
\]
qu'on appelera également (et peut-être à meilleur escient) la \uline{classe fondamentale locale} \uline{pour} $f$. On notera que pour la connaitre pour tout $A$, il suffit de la connaitre pour les $\underline{\mathrm{Z}}/n\underline{\mathrm{Z}}$ (du moins si on se borne aux $A$ qui sont annulés par un même entier $n\handwritten{\geq}1$ premier aux car. résiduelles de $S$): \deleted{La connaissance de (4.6) permet de reconstituer} plus précisément, la formation des classes $\handwritten{\phi}_f$ de (4.5) et de (4.6) est \uline{fonctoriel} en l'Anneau $A_S$ (voire, en $A_Y$).

Lorsqu'on dispose d'une ``formule de pureté''
\[
\deleted{\unsure{\underline{\mathrm{H}}^i(\mathrm{R}|\underline{\mathrm{H}}^i\mathrm{R}}}\ \mathrm{R}^i f^{\handwritten{!}}(A_Y)=0\quad\text{pour }i\handwritten{<}2(d-d'),
\tag*{(4.7)}
\]
la connaissance de (4.6) équivaut même à celle de (4.5), ce dernier $\handwritten{\phi}_f$ méritant le nom de \uline{classe fondamentale globale} de $X$ sur $Y$.

Nous allons nous intéresser par la suite au cas où $f:X\handwritten{\longrightarrow}Y$ est une \uline{immersion fermée}. On sait alors que les formules de pureté (4.7)
